\section{Problem Landscape}
Eisert and Preskill identify four gaps that must be bridged en route from NISQ demonstrations to fault-tolerant application-scale quantum computing. Our current tooling leaves several of these gaps unaddressed:
\begin{itemize}
  \item \textbf{Gap (iii)---from heuristics to mature algorithms.} High-level process models, typically Pauli-sum Hamiltonians or ZX-style diagrams, rarely flow into provably correct optimisation pipelines. Manual rewrites dominate, and the lack of certificates prevents downstream verification.
  \item \textbf{Gap (i)---from mitigation to detection and correction.} Error-mitigation strategies are often bolted on after circuit synthesis, blocking the feedback needed to inform code selection or hardware calibration.
  \item \textbf{Gap (iv)---from exploratory simulators to credible quantum advantage.} Extraction towards hardware-specific circuits seldom preserves the provenance information required for reproducible benchmarking or credible claims of advantage.
\end{itemize}

Partners emphasise that these blockers translate into delayed experiments, costly calibration cycles, and difficulty engaging with error-correcting code development. A practical yet rigorous compiler must therefore:
\begin{enumerate}
  \item Ingest hybrid coherent-stochastic models without fragmenting the reasoning formalism.
  \item Produce certified ZXW rewrite traces that scale to diagrams with hundreds of nodes and interface with verification tooling.
  \item Deliver circuit extraction pipelines that preserve traceability, expose error-mitigation hooks, and align with target gate sets and code architectures.
\end{enumerate}
